\section{Methodology}
\label{sec:method}

\subsection{The hypothesis, decomposed into testable clauses}
\label{sec:clauses}

We preregistered the hypothesis as a conjunction and fixed a refutation criterion for each
clause before fitting any model (\tabref{tab:clauses}). Clause (i) is tested by the pooled
association and, if null, by equivalence testing; clause (ii) by conversion to
$\Delta\mathrm{SUVR}$ per 10 \pp of N3; clause (iii) by within-study paired contrasts that must
additionally survive a measurement-reliability null model. A negative result on any clause is
reported as the finding.

\begin{table}[t]
    \centering
    \small
    \begin{tabular}{@{}cp{2.9cm}p{3.5cm}p{4.5cm}@{}}
        \toprule
        \textbf{Clause} & \textbf{Claim} & \textbf{Supported if} & \textbf{Refuted if} \\
        \midrule
        (i)   & N3 stage \% is associated with amyloid PET burden
              & pooled $r > 0$ with 95\% CI excluding 0
              & CI includes 0 \emph{and} excludes $r \geq 0.20$ (TOST) \\
        (ii)  & the association is dose-dependent and measurable in SUVR
              & $\Delta\mathrm{SUVR}$ per 10 \pp CI excludes 0 and is a non-trivial fraction of
                the A$-$/A$+$ separation
              & CI includes 0, \emph{or} magnitude $< \sim$5\% of the A$-$/A$+$ SUVR gap \\
        (iii) & it is \emph{stronger} than \tst and \se
              & paired $\Delta z > 0$ with CI excluding 0 \emph{and} the excess survives the
                reliability null
              & $\Delta z \leq 0$, \emph{or} any positive gap is fully explained by differential
                ICC \\
        \bottomrule
    \end{tabular}
    \caption{The hypothesis decomposed into three preregistered clauses with their refutation
    criteria, fixed before any model was fitted. The hypothesis is a conjunction: it is
    supported only if all three clauses hold.}
    \label{tab:clauses}
\end{table}


\subsection{Corpus and analysis set}
\label{sec:data}

The primary input is a curated table of 49 effect sizes hand-extracted from 17 studies across 12
independent cohort families. Every row carries a verified DOI or PMID, a verbatim provenance
quote from the source paper, a \texttt{cohort\_family} label used for dependency handling, and a
\emph{burden-aligned} effect signed so that \textbf{positive means worse sleep is associated
with more amyloid}. Positive therefore always supports the hypothesis, regardless of how the
original paper oriented its exposure.

The primary analysis set is the 40 usable effects whose outcome is amyloid PET burden
(cross-sectional deposition or change). The nine CSF amyloid-beta42 effects are excluded from
the primary model and analysed only as a cognitive-status-stratified secondary, because the sign
of soluble CSF amyloid-beta42 relative to deposited burden reverses between pre-deposition
samples \citep{varga2016} and post-deposition samples \citep{liguori2020}.

Three supporting datasets do the work the primary table cannot. Single-night ICC(2,1) values per
sleep metric come from \sleepedf \citep{kemp2000}, computed on 75 subjects with two consecutive
nights each, and feed the reliability null model. Subject-level N3 percentages from the same
source (153 recordings, 78 subjects) give the exposure dispersion for the dose conversion and
the feasibility check. Amyloid SUVR/Centiloid anchors come from \citet{jack2017}, and
\citet{chen2025} serves as an external benchmark for the pooled comparator estimates.

\subsection{Effect harmonisation}
\label{sec:harmonise}

Everything is pooled on the Fisher-$z$ scale \citep{fisher1915},
$z = \operatorname{artanh}(r)$ with $v = 1/(n-3)$. Kendall's $\tau$ is converted through
Greiner's relation $r = \sin(\pi\tau/2)$ \citep{kendall1949}. \Tabref{tab:conversion} lists the conversion rule applied to each input type and
how many amyloid-PET effects it covers.

\begin{table}[t]
    \centering
    \small
    \resizebox{\ifdim\width>\textwidth\textwidth\else\width\fi}{!}{%
\begin{tabular}{@{}llc@{}}
        \toprule
        \textbf{Input type} & \textbf{Conversion rule} & \textbf{PET effects} \\
        \midrule
        \texttt{pearson\_r}, \texttt{partial\_r}, \texttt{r\_from\_partial\_eta2}
            & used directly as $r$ (with $r = \sqrt{\eta_p^2}$ for a 1-df contrast) & 20 \\
        \texttt{kendall\_tau}
            & $r = \sin(\pi\tau/2)$ (Greiner's relation) & 1 \\
        \texttt{unstd\_beta\_*}
            & test-statistic recovery, \eqref{eq:recovery} & 12 \\
        \texttt{narrative\_null}
            & imputed $z = 0$, $v = 1/(n-3)$ & 7 \\
        \midrule
        \multicolumn{2}{@{}l}{\textbf{Total usable amyloid-PET effects}} & \textbf{40} \\
        \bottomrule
    \end{tabular}%
}
    \caption{Effect harmonisation rules, fixed in advance and applied uniformly. All effects are
    pooled as Fisher $z$ with $v = 1/(n-3)$. Test-statistic recovery raises the usable
    amyloid-PET set from 21 native correlations to 40 effects and is what brings the two largest
    cohorts onto the common metric.}
    \label{tab:conversion}
\end{table}


\para{Test-statistic recovery.} The step that makes this analysis possible is the recovery of a
partial correlation from a reported regression coefficient and its $p$-value
\citep{rosenthal1979}. For a two-sided $p$ on $\mathrm{df} = n - 1 - n_{\mathrm{cov}}$ residual
degrees of freedom,
\begin{equation}
t = \operatorname{sign}(\beta)\, F^{-1}\!\left(1 - \tfrac{p}{2};\, \mathrm{df}\right),
\qquad
r = \frac{t}{\sqrt{t^{2} + \mathrm{df}}}.
\label{eq:recovery}
\end{equation}
Raw betas are never pooled in their native units. Recovery brings the two largest cohorts
(\afour, $n = 4{,}417$; \aibl, $n = 189$) and the \mayo and \mars results onto the common
metric, raising the usable amyloid-PET set from 21 native correlations to 40 effects. The
conversion is exact where it can be checked: recovering $r$ from the $p$-value of a simple
correlation returns the original $r$ to within $10^{-8}$.

\para{Reported nulls.} Seven amyloid-PET effects are narrative nulls --- statements such as ``no
association with total sleep time or sleep efficiency'' with no number attached. We impute
$z = 0$ with $v = 1/(n-3)$. This is conservative in magnitude but anti-conservative in
precision, so it is carried as the primary rule \emph{and} removed in sensitivity analysis S1.
Dropping these rows instead would bias every pooled estimate away from zero.

\subsection{Dependency handling and estimation}
\label{sec:estimation}

Berkeley alone contributes 13 of 49 effects, so treating publications as independent would
inflate precision. The primary estimator first aggregates to one synthetic effect per cohort
family per stratum using the correlated-effects variance of \citet{borenstein2009},
\begin{equation}
\bar{z} = \frac{1}{m}\sum_{i=1}^{m} z_i,
\qquad
\bar{v} = \frac{1}{m^{2}}\left(\sum_{i=1}^{m} v_i + \sum_{i \neq j} \rho \sqrt{v_i v_j}\right),
\qquad \rho = 0.6,
\label{eq:borenstein}
\end{equation}
and then fits REML random effects \citep{viechtbauer2005} with
Knapp--Hartung--Sidik--Jonkman standard errors on a $t_{k-1}$ reference distribution
\citep{knapp2003}, the recommended small-$k$ adjustment. Naive per-effect pooling is reported
only to quantify how much dependency handling matters, and cluster-robust variance estimation
\citep{hedges2010} is reported where a stratum has at least four cohort families.

\subsection{The six analyses}
\label{sec:experiments}

\para{E1: stratified pooling.} Pooled estimates for \stage, \spectral, \tst, \se, \waso and
other stages, with $\tau^2$, $Q$, $I^2$ and a 95\% prediction interval \citep{higgins2009}.
Because $k \leq 6$ cohort families per stratum, every primary $p$-value is accompanied by an
exact sign-flip permutation $p$-value enumerating all $2^k$ assignments.

\para{E2: specificity.} A naive between-stratum comparison is invalid because the strata share
cohorts. The primary test is therefore a within-study paired difference: for every cohort family
reporting both exposures in the same participants,
$\Delta z = z_{\mathrm{N3}} - z_{\mathrm{comparator}}$ with
$\operatorname{Var}(\Delta z) = v_a + v_b - 2\rho_w\sqrt{v_a v_b}$ and $\rho_w = 0.5$. Three
preregistered confirmatory contrasts (stage vs.\ \tst, stage vs.\ \se, stage vs.\ \spectral)
carry a Holm correction \citep{holm1979}.

\para{E3: measurement-reliability null model.} Classical attenuation gives
$r_{\mathrm{obs}} = r_{\mathrm{true}}\sqrt{\mathrm{ICC}_{\mathrm{exposure}}}
\sqrt{\mathrm{ICC}_{\mathrm{outcome}}}$. In a within-outcome comparison the outcome term
cancels, so if two exposures had \emph{identical} true correlations with amyloid their observed
correlations would still stand in the ratio
$\sqrt{\mathrm{ICC}_a / \mathrm{ICC}_b}$. We compute that predicted ratio from the empirical
single-night ICCs and compare it to the observed ratio, and we also report each stratum
disattenuated for its own unreliability.

\para{E4: dose conversion and feasibility.} We convert the pooled correlation into the units the
hypothesis states,
\begin{equation}
\Delta\mathrm{amyloid} \text{ per } 10\,\mathrm{pp\ N3}
= r_{\mathrm{pooled}} \times \frac{10}{\mathrm{SD}_{\mathrm{N3\%}}}
\times \mathrm{SD}_{\mathrm{amyloid}},
\label{eq:dose}
\end{equation}
propagating all three inputs by Monte Carlo (10{,}000 draws): $r$ via a $t_{k-1}$ draw on the
Fisher-$z$ scale, $\mathrm{SD}_{\mathrm{N3\%}}$ via the $\chi^2$ sampling distribution of an SD,
and $\mathrm{SD}_{\mathrm{amyloid}}$ via a $\mathrm{Beta}(30,70)$ prior on amyloid-positive
prevalence applied to the medians and IQRs of \citet{jack2017}. We then ask whether a
10-percentage-point reduction in N3 is an available exposure contrast at all, using the observed
N3 distribution in \sleepedf subjects aged 60 and over.

\para{E5: robustness and bias.} Seventeen preregistered specifications: dropping imputed nulls
(S1), dropping $p$-recovered betas (S2), both (S1+S2), alternative residual-df rules (S3),
$\rho \in \{0.3, 0.8, 1.0\}$ (S4), cognitively normal samples only (S5), PSG-measured exposures
only (S6), leave-one-cohort-family-out (S7), fixed effects (S8), and the CSF secondary (S10).
Small-study bias by Egger regression \citep{egger1997} and Duval--Tweedie trim-and-fill
\citep{duval2000}, labelled uninterpretable below $k = 5$ and $k = 10$ respectively.

\para{E6: precision.} Minimum detectable effect at 80\% power per stratum, power at
$r \in \{0.20, 0.30, 0.45\}$, and equivalence testing \citep{lakens2017} of the \stage stratum
against bounds $r = \pm 0.10, 0.15, 0.20, 0.30, 0.45$. An honest negative result must state what
it could and could not have detected.

\subsection{Implementation and reproducibility}
\label{sec:impl}

Python 3.12.8 with numpy 2.5.2, pandas 3.0.5, scipy 1.18.1, statsmodels 0.15.0 and pymare
0.0.12, in an isolated \texttt{uv} environment with pinned dependencies. Seed 42 throughout;
the only stochastic step is the E4 Monte Carlo. The full pipeline runs in about 29 seconds on
CPU. GPUs were available but not used: this is a study-level meta-analysis whose entire compute
cost is trivial, so acceleration would add complexity with no benefit.

The meta-analytic primitives (REML $\tau^2$, HKSJ, cluster-robust variance, exact permutation,
trim-and-fill, TOST) are implemented directly from their primary formulations rather than taken
from a package, because no available Python package provides HKSJ, cluster-robust variance and
exact permutation together. They are cross-checked against PyMARE, an independently maintained
meta-analysis package: REML point estimates agree to within $4.4 \times 10^{-9}$ on all four
strata. Validation re-runs the entire pipeline in a scratch copy and confirms all eight result
files are bit-identical; 39 of 39 validation checks pass.
